% !PRINTED 372
\weberpage{372}

The substitution (5) is said to be \gesperrt{composed} of (1) and (4).

The need for an abbreviated notation for these substitutions
makes itself felt. Since it is often not the variables that matter,
but only the substitution numbers $\alpha, \beta, \gamma, \delta$,
one denotes the whole substitution (1) by the substitution numbers,
or also merely by a single letter
\weq{9}{S = \begin{pmatrix} \alpha, & \beta \\ \gamma, & \delta \end{pmatrix},}

and then, when it is necessary, writes the equation (1) thus:
\[
y = S(x) = \begin{pmatrix} \alpha, & \beta \\ \gamma, & \delta \end{pmatrix}(x).
\]

The composition of two substitutions is denoted by setting
the symbols side by side, in which, however, attention must be
paid to the \gesperrt{order}; thus
\weq{10}{\begin{pmatrix} \alpha, & \beta \\ \gamma, & \delta \end{pmatrix} \begin{pmatrix} \alpha', & \beta' \\ \gamma', & \delta' \end{pmatrix} = \begin{pmatrix} \alpha'', & \beta'' \\ \gamma'', & \delta'' \end{pmatrix},}

or
\weq{11}{S S' = S''.}

The formulas (6) contain the rule according to which a
composite substitution is to be formed. One can derive it from
the multiplication rule for determinants, but must observe that
in the first factor the summation is to be taken by rows, in the
second by columns, just as the formulas (6) show.

Just as one composes two substitutions, one can also form
the composition of several:
\[
S S_1 S_2 S_3 \ldots
\]

In forming the composite substitution the components may
not, in general at least, be interchanged; but one may at will
combine two adjacent ones into one, then again two, and so
on; that is, the \gesperrt{commutative} law does not hold, but the
\gesperrt{associative} law does.

This follows immediately from the fact that, if $x$ is expressed
in terms of $x_1$, $x_1$ in terms of $x_2$, and $x_2$ in terms of $x_3$,
then the expression of $x$ in terms of $x_3$ can be found either in
this way, that one first expresses $x_1$ in $x$ by means of $x_2$ and
then $x_2$ by means of $x_3$, or in this way, that
%CONT>>
